% Proposed replacement for the Introduction in the v2 manuscript.
% Uses existing citation keys, plus castellano2012 in refs_additions.bib.
\section{Introduction}
\label{sec:intro}

A transition into an absorbing state eliminates fluctuations. When two absorbing
configurations are related by a $\mathbb{Z}_2$ symmetry, it can also select one of two
equivalent ordered states. In two dimensions, this combination is closely associated with
generalized voter (GV) criticality and its unusually slow coarsening. We study a binary
lattice model in which the active phase instead meets the two absorbing states at a
discontinuous transition. The evidence is a finite interface density at the transition,
reversing fronts, metastable decay, and nucleation, measured in complementary geometries.

The two-dimensional picture developed through several distinct results. Dornic \emph{et al.}
identified voter coarsening as a universality class for transitions driven by interfacial
noise between dynamically equivalent absorbing states~\cite{dornic2001}. Its prototype is the
voter model, whose interface density decays as $1/\ln t$ in two
dimensions~\cite{clifford1973,holley1975,cox1986,krapivsky1992}. Droz, Ferreira, and Lipowski
then showed that extending the interaction range can split a direct transition into Ising
symmetry breaking followed by directed-percolation absorption~\cite{droz2003}. The Langevin
description of Al Hammal \emph{et al.} unified these scenarios: although its deterministic
drift permits bistability, their two-dimensional simulations gave a GV transition, whereas
they reported fully first-order behavior in three dimensions~\cite{alhammal2005}. Subsequent
microscopic studies recovered GV and Ising--DP transitions~\cite{vazquez2008,castellano2009,azizi2020}.
Rodrigues \emph{et al.}, however, reported a discontinuous line for a triangular-lattice
model~\cite{rodrigues2015}. Their analysis does not establish the finite activity jump and
phase-coexistence properties tested here, leaving a reason to examine these observables
directly.

The distinction matters because a magnetization jump alone does not establish a
discontinuous loss of activity: it also occurs at the GV transition, where the interface
density vanishes continuously~\cite{dornic2001,alhammal2005}. Here a discontinuous transition
means a nonzero limiting activity on the active side, supported by coexistence
and nucleation tests. Such behavior is established in two-dimensional reaction
models~\cite{ziff1986,bidaux1989,lubeck2006}, in models with three symmetric absorbing
states~\cite{lipowski2002,azizi2018}, and in cyclic ecological models whose two absorbing
alliances contain many configurations and retain internal motion~\cite{szabo2004,guisoni2013}.
The latter provide close precedents for front competition and active layers between ordered
phases. The microscopic setting here is different: each site has two states, and there are
exactly two absorbing configurations.

Our model combines the self-complementary Day \& Night cellular-automaton rule with voter
copying and a local minority flip. It is a short-range nonlinear voter model: the flip
probability depends only on the number of disagreeing Moore neighbors. This construction
connects it to the nonlinear voter models of Cox and Durrett~\cite{cox1991} and the
probabilistic automata studied by Molofsky \emph{et al.}~\cite{molofsky1999}. A positive voter
weight removes all frozen configurations except the two uniform states. At the
representative weight $p=1/2$, we locate the transition from planar-front motion and test
the bulk and droplet dynamics around that location. We compare different front
orientations and place a large droplet of each phase inside the other, which checks for
the orientation-dependent coexistence possible in nonequilibrium
models~\cite{wang2015,liu2018}. Bulk measurements then resolve a dense active phase and its
metastable decay. Walls between the two absorbing phases acquire an active layer whose
thickness increases toward the transition.

We also follow the invasion boundary toward the voter point. Coalescing random walks
provide a finite-time closure for the reaction term of a weak voter perturbation, inspired
by the rigorous theory in $d\ge3$~\cite{cox2013}. Within this closure the transition scale is
$q^*\propto\alpha(1-p)$, with an amplitude close to the measured one. Contributions from
five or more distinct ancestral lineages can support bistability. In two dimensions the
matching time remains undetermined, so this construction is an approximate account of the
boundary and a guide to the endpoint problem.

The five-cell automaton of Molofsky \emph{et al.} supplies a useful comparison. Their
mean-field map has a regime with three stable fixed points, but their spatial simulations
did not find the corresponding dependence on initial density~\cite{molofsky1999}.
Our closure approximately locates its invasion boundary, while direct simulations reveal
strong suppression of small seeds near that boundary. They do not resolve the order of
the transition or the metastable active branch predicted by the closure. This comparison
separates the successful boundary estimate from the stronger claim of bistability.

Section~\ref{sec:model} defines the dynamics and observables.
Section~\ref{sec:flagship} presents the coexistence tests at $p=1/2$.
Section~\ref{sec:line} studies the parameter dependence and sensitivity to update timing,
and Sec.~\ref{sec:theory} develops the finite-time closure.
The comparison with the five-cell rule is given in Appendix~\ref{app:molofsky}.
